\begin{theorem}[B.4.5 (Proved): the obstruction on normalization components]
Let $D$ be a compact reduced complex space with normal-crossings
singularities, allowing self-intersections of global components.
Suppose that every connected component of its smooth normalization
satisfies the $\partial\bar\partial$-lemma. Then the natural map
\[
 H^1(D;\C)\longrightarrow H^1(D,\OO_D)
\]
is surjective.

Consequently every smooth-total-space reduced normal-crossings
filling of the specified six-sphere family, after any finite base
change, has a normalization component which fails the
$\partial\bar\partial$-lemma. In particular its normalization
cannot consist entirely of compact K\"ahler or Moishezon manifolds.
This conclusion includes normal crossings which are not simple.
\end{theorem}
\begin{proof}
We give a degree-one argument directly on $D$; no nearby
Hodge-de Rham degeneration theorem is needed. Let
$D^{[p]}$ be the smooth compact branch space whose points are
ordered $(p+1)$-tuples of distinct local branches through one
point of $D$. Equivalently, we normalize the loci where these
branches meet, retaining their ordering. These are distinct
\emph{local branches}, even when they belong to one global
irreducible component. The finite holomorphic maps
$\pi_p:D^{[p]}\to D$ and their face maps are specified by
forgetting one branch. The symmetric group acts by reordering;
we take the alternating summand, with the sign character, in
each degree. Averaging over this finite group is exact over
$\C$. Write $\delta$ for the alternating sum of the face
pullbacks.

Both augmented complexes
\[
 \C_D\longrightarrow
 \bigl((\pi_p)_*\C_{D^{[p]}}\bigr)_{p\geq0}^{\mathrm{alt}},
 \qquad
 \OO_D\longrightarrow
 \bigl((\pi_p)_*\OO_{D^{[p]}}\bigr)_{p\geq0}^{\mathrm{alt}}
\]
are resolutions. This can be checked locally without a global
ordering. In a chart $z_1\cdots z_k=0$, the constant complex is
the augmented cochain complex of the simplex on the $k$ branches.
For the holomorphic complex, fix a monomial in the convergent
power series ring. It survives on exactly the intersections
indexed by subsets of those $i$ for which its $z_i$-exponent is
zero. If there are no such indices it already vanishes on $D$;
otherwise its coefficient complex is again an augmented simplex
complex. The usual contraction is a finite combination of
coordinate restrictions, so preserves convergence. This proves
exactness for analytic functions, not just for formal series.
The descriptions agree under permutations of the local branches.
The finite simplex and permutation identities for the at most
four branches in our threefold application are checked exactly
in \path{code/b_nc_degree_one.py}.

Resolve the holomorphic terms by smooth Dolbeault forms and
the constant terms by smooth de Rham forms on the branch spaces.
These compute the respective hypercohomologies. Indeed finite
pushforward is exact here, and partitions of unity from local
ambient smooth functions pull back to make the resulting sheaves
fine. The ordinary de Rham lemma and Dolbeault lemma apply on
each smooth branch space. Projection to the $(0,\bullet)$
part defines the natural comparison between the two total
complexes and induces $H^1(D;\C)\to H^1(D,\OO_D)$.

Represent a coherent degree-one class by a total cocycle
$(\alpha,f)$, where $\alpha$ is a smooth $(0,1)$-form on
$D^{[0]}$ and $f$ is an alternating smooth function on
$D^{[1]}$. Choose the total-complex signs so that its equations
are
\[
 \bar\partial\alpha=0,\qquad
 \delta\alpha=\bar\partial f,\qquad \delta f=0.
\]
On a normalization component, $\partial\alpha$ is of type
$(1,1)$, is $\partial$-exact and is $\bar\partial$-closed.
The $\partial\bar\partial$-lemma gives a smooth function $h$
with $\partial\alpha=\partial\bar\partial h$.
Do this on every component, and put
\[
 \alpha'=\alpha-\bar\partial h,\qquad
 f'=f-\delta h.
\]
The pair represents the same coherent class and still satisfies
the displayed cocycle equations; now $d\alpha'=0$.
Holomorphic face maps commute with both differential operators,
so
\[
 \partial\bar\partial f'
       =\partial\delta\alpha'=\delta\partial\alpha'=0.
\]
A smooth complex-valued pluriharmonic function on a compact
connected complex manifold is constant: apply the maximum
principle to its real and imaginary parts in local complex
coordinates. Therefore $f'$ is constant on each connected
component of $D^{[1]}$. It follows that
$df'=0$ and $\delta\alpha'=0$.

Thus the same pair $(\alpha',f')$, with $\alpha'$ viewed as
a complex-valued smooth one-form, is a cocycle of the total
de Rham branch complex. Its image is the original coherent
class. This proves surjectivity. Notice that no
$\partial\bar\partial$ hypothesis on $D^{[1]}$ or higher
branch spaces was used; compactness suffices for the only
function argument there.

For the asserted application, B.4.2 and B.4.3 give
$H^1(D;\C)=0$ and $h^1(D,\OO_D)\geq1$.
Surjectivity is impossible, so some normalization component
fails the lemma. Compact K\"ahler and smooth Moishezon
manifolds satisfy it, as in the primary comparison checked
in B.4.3. The same argument applies after every finite base
change, since its smooth fibres are still the specified
six-spheres.
\end{proof}

This strengthens the simple-normal-crossings obstruction in
B.4.3 by using only degree one. It also avoids an unverified
extension of an algebraic Du Bois theorem to arbitrary analytic
spaces. Searches for normal-crossings cohomology with K\"ahler
normalization, the normalization Dolbeault complex, and
degree-one surjectivity located the broader Hodge-theoretic
setting discussed by \citet{Chen2026HodgeRoot} and the proper
algebraic statement of \citet{GrafKovacs2014Root}; we do not
claim priority for this elementary consequence. The proof above
states and verifies the particular analytic comparison required
here. It still does not prove that all possible analytic
normalization components satisfy its hypothesis.
